% Paper 6 (GRPO-Registry) appendix: artifact reference
\section{Artifact Reference}
\label{sec:p6-appendix}

\subsection{Layout and Validation}
\label{sec:p6-appendix-layout}
The artifact ships in the benchmark repository:

\begin{lstlisting}[basicstyle=\ttfamily\scriptsize, frame=single,
  caption={Artifact layout.}, label={lst:p6-layout}]
registry/
  schema.json          # JSON Schema, draft 2020-12; record types:
                       #   stack | variant_delta
  entries/             # 28 stack + 18 variant_delta records; 2 amendment sidecars
    trl_grpo_qwen3-8b_gsm8k.json        verl_grpo_qwen3-8b_gsm8k.json
    openrlhf_grpo_qwen3-8b_gsm8k.json   tinker_grpo_qwen3-8b_gsm8k.json
    colab-open_grpo_e3.json             colab-open_drgrpo_e3.json
    colab-open_dapo_e3.json             colab-open_grpo-adaptiveg_e3.json
    tinker_grpo_qwen3.5-4b_gsm8k.json   tinker_dapo_qwen3.5-4b_gsm8k.json
    tinker_drgrpo_qwen3.5-4b_gsm8k.json tinker_gspo_qwen3.5-4b_gsm8k.json
    delta_dapo.json  delta_drgrpo.json  delta_gspo.json
  query.py             # reference CLI (stdlib only)
  README.md
\end{lstlisting}

Every entry validates against the schema:

\begin{lstlisting}[basicstyle=\ttfamily\scriptsize, frame=single]
python3 -c "import json, glob, jsonschema; \
  s = json.load(open('registry/schema.json')); \
  [jsonschema.validate(json.load(open(p)), s) \
   for p in glob.glob('registry/entries/*.json')]"
\end{lstlisting}

\subsection{A Complete Variant-Delta Record}
\label{sec:p6-appendix-delta}
The DAPO record, verified against the source paper~\citep{yu2025dapo}:

\begin{lstlisting}[basicstyle=\ttfamily\scriptsize, frame=single,
  caption={\texttt{entries/delta_dapo.json} (whitespace compressed).},
  label={lst:p6-delta-dapo}]
{ "record_type": "variant_delta", "schema_version": "0.1.0",
  "id": "delta_dapo", "name": "DAPO", "base": "grpo",
  "citation": { "bibkey": "yu2025dapo", "arxiv": "2503.14476",
    "title": "DAPO: An Open-Source LLM Reinforcement Learning
              System at Scale" },
  "deltas": [
   {"component": "clip_higher",
    "field": "loss_form.clip_eps_high",
    "change": "asymmetric (decoupled) clip: eps_low=0.2,
               eps_high=0.28"},
   {"component": "dynamic_sampling",
    "field": "sampling.dynamic_sampling",
    "change": "filter prompts whose groups have zero reward
               variance and resample until the batch is full"},
   {"component": "token_level_loss",
    "field": "loss_form.aggregation",
    "change": "token-level policy-gradient loss (mean over all
               tokens in the batch rather than per-sequence
               averaging)"},
   {"component": "overlong_reward_shaping",
    "field": "reward.overlong_shaping",
    "change": "soft length-aware penalty for truncated/overlong
               completions"},
   {"component": "kl_removed", "field": "reference_kl",
    "change": "KL regularization term removed from the objective"}],
  "notes": "Four named techniques plus removal of the KL term,
            per the source paper." }
\end{lstlisting}

\subsection{Badge Output on the Seed Entries}
\label{sec:p6-appendix-badge}
The current per-entry badge and specification coverage are regenerated from
\texttt{registry/query.py} in Table~\ref{tab:p6-entries}. The earlier hard-coded
Listing 7 used an older scoring schema and is removed; its scores must not be
compared with current badges. Per-item fractions are printed by
\texttt{python3 registry/query.py badge}; the denominator is the canonical
leaf set in that implementation, not an arbitrarily chosen number of
populated JSON keys. Outcomes and CI-method declarations are not scored.
\endinput
// ... historical badge output below omitted from manuscript ...
